% LSPV-5 / GLOBALB2GATE1 / condensed durable manuscript / November 18 2026 01:10 PST
\documentclass[11pt]{article}
\usepackage{amsmath,amssymb,amsthm,hyperref}
\newtheorem{theorem}{Theorem}
\title{GLOBALB2GATE1: The Global Two-Output Gate Is coNP-Complete}
\author{LSPV-5 / source-grounded certificate research}
\date{18 November 2026}
\begin{document}\maketitle
\noindent\textbf{STATUS.} PROVED/EXACT for explicit finite-source reductions below, independent mathematical review OPEN. Five exact test checkpoints COMPUTATIONAL. Universal PFCERT, nonalgebrization, and P versus NP OPEN. This is a standalone original-source result, not a lower bound for unrestricted SAT.
\section{IDEAS PROPOSED BEFORE LITERATURE}
Original frozen proposals SHA256 6c55d5795fb5dabd7165ff7d8a29d7f10c9dd8e1920bdb6250de8cd97a97ca1e, retained in full local archive.
\par\textbf{CANDIDATE A.} Given a fixed polynomial source-grounded proof checker, decide if at most two distinct outputs are accepted across ALL proofs. Missing canonical finite verifier and exact cost; falsifier: a polynomial solution to a coNP-complete gate. 
\par\textbf{CANDIDATE B.} Binary-horizon exact endpoint of a succinct Boolean successor circuit is PSPACE-complete despite polynomial-time single-step evaluation. Missing faithful configuration/TM proof and horizon accounting; falsifier: a uniform polynomial NP-verifiable exact singleton endpoint certificate would imply PSPACE subset NP.
\par\textbf{CANDIDATE C.} Exponentially many individually singleton sparse descriptors may collectively accept every possible output. Missing exact proof-string projection; falsifier: enumerate singleton descriptors.
\section{LITERATURE CHECK AND REVISIONS}
Primary Hirahara (STOC 2020/ECCC TR20-050, Theorem 3.9, Section 9) proves P-immunity of superpolynomial-time-bounded shortest-program K^t randomness using retained n-O(log n) prefix and polynomial suffix search, distinct from Levin Kt. Primary Baker--Gill--Solovay, SIAM J. Comput. 4 (1975), construct a recursive A with P^A=NP^A using deterministic/nondeterministic query machines with shared oracle-query semantics. Galperin--Wigderson, Information and Control 56 (1983), study succinct graph encodings; a 2017 Information and Computation paper (doi:10.1016/j.ic.2017.07.008) proves related Boolean dynamical-system PSPACE results. Neither is substituted for the direct reduction below. A--C survive in restricted form. A polynomial local check does not imply a polynomial global bound across all proofs.
\section{CONSTRUCTION OR PROOF ATTEMPT}
Fix a canonical complete CNF source F on m variables, a bit length n>=max(m,2), and a fixed deterministic polynomial source interpreter V. Fix program tag p=empty and three distinct n-bit outputs x0=0^n, x1=10^{n-1}, x2=010^{n-2}. Proofs are tags 0,1,2 and m-bit assignments a. V accepts x0 under tag 0 and a=0^m, accepts x1 under tag 1 and a=0^m, and accepts x2 under tag 2 precisely if a satisfies every clause of F; reject all else. Every individual acceptance checks finite source semantics in polynomial time and names exactly one output; conservative sequential-tape time O((|F|+n)^3) and writable space O(|F|+n).
\begin{theorem}[Global B2 coNP-completeness]
For this fixed V, the language B2(F) saying that at most two distinct x have some accepted proof is coNP-complete under polynomial many-one reductions. For ANY fixed polynomial verifier and polynomially bounded x and proof lengths, global at-most-two membership lies in coNP.
\end{theorem}
\noindent\emph{Proof.} Failure has three different polynomial strings x_i and their three accepted polynomial proofs. In the constructed fixed V, accepted image equals {x0,x1} iff F is UNSAT, and {x0,x1,x2} otherwise. Thus F maps to B2(F) iff F is UNSAT; source construction is polynomial. \qed
\begin{theorem}[CoNP-complete invariant coverage despite trivial dynamics]
Let a configuration be (flag,u) in {0,1}^{m+1}, transition be the identity map, initial state (0,0^m), all states terminal with output flag, and proposed invariant I_F(flag,u)=[flag=0 AND u=0^m] OR [flag=1 AND F(u)]. The verification of initiality, consecution and terminal coverage with allowed envelope {0} is coNP-complete. Initiality and consecution always hold; coverage holds exactly when F is UNSAT.
\end{theorem}
\noindent\emph{Proof.} The invariant contains the initial state and is inductive under identity; a bad terminal configuration in the invariant exists exactly for a satisfying F assignment. This is polynomially verifiable and gives the UNSAT reduction. \qed
\par\textbf{Scope.} The actual reachable set is the singleton initial state, and a different trivial invariant proves its sole endpoint. The theorem is about AUDITING a supplied weak invariant, not proving no succinct invariant exists.
\begin{theorem}[Singleton-union counterexample]
Proof d in {0,1}^n can name the source-certifiable singleton descriptor E_d={d}. Each individual descriptor has computable cardinality one, yet the existential projection of all accepted proofs has exactly 2^n different outputs.
\end{theorem}
\noindent\emph{Proof.} For each x choose d=x. \qed
\begin{theorem}[Binary-horizon circuit iteration]
Given an explicit Boolean gate DAG C:{0,1}^m to {0,1}^m, start s, binary T and output bit index i, the language [(C^T(s))_i=1] is PSPACE-complete.
\end{theorem}
\noindent\emph{Proof.} Store m-bit current state, log(T+1)-bit counter and polynomial gate-evaluation scratch for PSPACE membership. For a fixed deterministic polynomial-space decider M for any PSPACE language, compile its authentic next-configuration transition into a polynomial gate DAG C_y, make accepting/rejecting halts absorbing, and use T=2^m (m+1 bits); a halting deterministic computation cannot repeat before halting and there are at most 2^m configurations. The persistent accept bit equals M(y). \qed
\par A polynomial-length, NP-checkable, GLOBALLY SINGLETON complete endpoint certificate for ALL such circuits would imply PSPACE subset NP. This implication is NOT valid for mere two-valued Boolean endpoint lists; both outputs can be admitted. Single-step/line access does not provide total explicit expansion or overcome exponentially large configurations in unrestricted machines.
\par\textbf{Matched-oracle PFCERT diagnostic.} For every oracle O use the SAME O in efficiently universal machine U^O, polynomial oracle verifier V^O, classes P^O and NP^O and randomness K_U^{t,O}. If an infinite P^O subset of R^{t,O}={x:K_U^{t,O}(x)>=|x|-1} existed for a superpolynomial t eventually dominating polynomial oracle search, store n-c log n prefix bits of an accepted string, self-delimit n with O(log n) bits and find a suffix by n^c membership queries to the SAME oracle decider. This would produce an accepted random string in <n-1 description bits within t(n), contradiction (relativized Hirahara prefix/suffix argument). If V^O covered every n-bit endpoint from all programs of length <=n-2, with at most two accepted outputs globally per (n,p) across ALL proofs, then E_n^O is in NP^O and |E_n^O|<=2(2^{n-1}-1)=2^n-2. Its infinite complement is contained in R^{t,O}. For a BGS collapsing oracle A with P^A=NP^A, the complement would belong to P^A and contradict immunity. Thus this UNIVERSAL two-output PFCERT schema cannot relativize to A. No ordinary existence theorem, nonalgebrization result or cryptographic hardness follows.
\par\textbf{Controls/resource ledger.} DEAD/COPY/DELAY/INTERP/RECOMP source variants, affine fast-forward and short postdisclosure hash-prefix search are legal shortcuts. Local V verification, source reading, parsing, proof and output handling are polynomial; the exponential global accepted-proof projection is NOT thereby polynomial. DTISP forbidden-region and general lower-bound gates OPEN; deterministic configuration counting is upper-only; rank/DISJ sanity PASS; universality and nonalgebrization OPEN.
\section{TESTS AND OUTCOME}
Five independent seeded Python checkpoint batches: B2 checker 768 cases / 1536 assertions; terminal coverage 736 / 2208; singleton union 12 / 8202; affine/finite iteration 909 / 909; source wrappers and disclosed prefixes 563 / 2803. Total 2988 finite cases, 15658 exact assertions, zero final failures. Seeds 2026111801--2026111805; source and raw JSON files in local evidence bundle. All results COMPUTATIONAL and not independent mathematical proof review.
\par\textbf{Cross-track ledger.} TMR-1 pinned-formalization and NP transfer OPEN; TMR-2 general necessary-work obstruction OPEN; TMR-3 complete blank-start checkpoint/evaluator OPEN; TMR-4 unrestricted cost-sound compiler transfer OPEN; TMR-5 full independent separation chain OPEN; LSPV-5 universal source-grounded nonrelativizing PFCERT existence OPEN. Peer archival reports inspected lack the four exact headings, though candidates/constructions are present in alternative sections, so do not infer absence of work. Historical main.tex is a 38-byte placeholder, not touched; current November 15 governance forbids root restoration. No schedules or peer-owned files changed.
\par\textbf{ONE next target.} Construct a uniform original-source invariant normal form for a specified nontrivial succinct transition class with syntactically GLOBAL sparsity over all accepted proofs, separately proving initiality, consecution, terminal coverage, state width, discovery and certificate costs, against PSPACE iteration and matched-oracle controls.
\begin{thebibliography}{9}
\bibitem{hirahara} Hirahara, STOC 2020/ECCC TR20-050, https://eccc.weizmann.ac.il/report/2020/050/.
\bibitem{bgs} Baker, Gill and Solovay, SIAM J. Comput. 4(4), 431--442 (1975), doi:10.1137/0204037.
\bibitem{gw} Galperin and Wigderson, Information and Control 56(3), 183--198 (1983).
\end{thebibliography}
\end{document}
% FINAL_VERIFICATION_SENTINEL::LSPV5::GLOBALB2GATE1::2026-11-18_0110::P_VS_NP_OPEN
